\section*{Chapter \ref{ch:g12:aggregate-demand-and-supply} --- Aggregate Demand and Supply}
% figdata: figdata/grade-12/08-aggregate-demand-and-supply/out/ (gdp, gap, recessions, adas_key); Kestria invented.

\begin{solution}{exo:g12:aggregate-demand-and-supply:1}
2001 (March to November), 2007--09 (December 2007 to June 2009) and 2020 (February to April).
\end{solution}

\begin{solution}{exo:g12:aggregate-demand-and-supply:2}
November 1973 to March 1975: 16 months. December 2007 to June 2009: 18 months.
\end{solution}

\begin{solution}{exo:g12:aggregate-demand-and-supply:3}
$(95 - 100)/100 = -5\,\%$.
\end{solution}

\begin{solution}{exo:g12:aggregate-demand-and-supply:4}
$(10.6 - 10)/10 = +6\,\%$: a boom, output above potential.
\end{solution}

\begin{solution}{exo:g12:aggregate-demand-and-supply:5}
Crash: demand. Harvest failure: supply. Orders from abroad: demand. Oil price: supply. Lockdown: both.
\end{solution}

\begin{solution}{exo:g12:aggregate-demand-and-supply:6}
$300 - 2P = 2P - 100$: $P = 100$, $Y = 100$.
\end{solution}

\begin{solution}{exo:g12:aggregate-demand-and-supply:7}
$260 - 2P = 2P - 100$: $P = 90$, $Y = 80$.
\end{solution}

\begin{solution}{exo:g12:aggregate-demand-and-supply:8}
Gap $-20\,\%$. In the long run output returns to 100; on the new demand curve, $100 = 260 - 2P$ gives $P = 80$.
\end{solution}

\begin{solution}{exo:g12:aggregate-demand-and-supply:9}
Long run: the price level rises by 10\,\%, output is unchanged. Short run: with prices sticky, the extra money raises spending and output rises; prices rise only as they
adjust, and output goes back to potential.
\end{solution}

\begin{solution}{exo:g12:aggregate-demand-and-supply:10}
$(22.731/3.431)^{1/64} - 1 \approx 3.0\,\%$ a year.
\end{solution}

\begin{solution}{exo:g12:aggregate-demand-and-supply:11}
1982, at $-5.2\,\%$. 1975: $-3.8\,\%$; 2009: $-2.9\,\%$; 2020: $-3.3\,\%$.
\end{solution}

\begin{solution}{exo:g12:aggregate-demand-and-supply:12}
With flexible prices, output stays at potential: a fall in demand only lowers prices, here by enough to keep real spending unchanged. In 2009 output fell by 2.6\,\% while prices
hardly moved: that contradicts flexible prices in the short run. The long-run part stands: output came back, and over decades money and prices still move together.
\end{solution}

\begin{solution}{exo:g12:aggregate-demand-and-supply:13}
A supply shock alone: prices up, output down. A demand recovery alone: prices and output up. Output rose by 6.2\,\%, so the demand recovery had the larger effect on output;
both shifts raised prices, so prices rose fast while output also rose.
\end{solution}

\begin{solution}{exo:g12:aggregate-demand-and-supply:14}
Year 1: $Y = 80$, $P = 90$. Year 2: gap $-10$, $Y = 90$, $P = (260 - 90)/2 = 85$. Year 3: gap $-5$, $Y = 95$, $P = 82.5$. The price level approaches 80.
\end{solution}

\begin{solution}{exo:g12:aggregate-demand-and-supply:15}
A straight line assumes \omterm{def:g12:aggregate-demand-and-supply:potential-output}{potential output} grew at one rate for 64 years. Growth was faster in the 1960s and slower after 2008, so the line lies above the data at both ends and
below it around 2000: the ``gaps'' it gives are mostly changes in trend growth, not slumps and booms. A smooth trend follows the data more closely, but it is still a
statistical line, not an estimate of the output possible at normal use of labour and capital: a long slump drags it down, so it understates the gap, and at the end of the
sample it depends on data that have not yet arrived.
\end{solution}

\begin{solution}{pb:g12:aggregate-demand-and-supply:1}
\textbf{1.} $P = Y$ and $Y = 200 - P$: $Y = 100$, $P = 100$.

\textbf{2.} Zero.

\textbf{3.} A vertical line at $Y = 100$.

\textbf{4.} With a given quantity of money, higher prices leave room for less real spending; Kestria's goods also become dearer than foreign ones.

\textbf{5.} With wages and other costs fixed for a time, a higher price level makes more production profitable.

\textbf{6.} $Y = 200 - (Y + 20)$: $Y = 90$, $P = 110$.

\textbf{7.} $-10\,\%$.

\textbf{8.} By 10\,\%. Prices up and output down: a supply shock.

\textbf{9.} On short-run supply with $Y = 100$: $P = 120$.

\textbf{10.} With output below potential, wages and other costs fall back, so $z$ shrinks and short-run supply shifts down.

\textbf{11.} 20, 12, 7.2, 4.32, 2.592.

\textbf{12.} $Y = 200 - P = 200 - Y - z$, so $2Y = 200 - z$ and $Y = 100 - z/2$.

\textbf{13.} Output 90, 94, 96.4, 97.84, 98.704; gaps $-10$, $-6$, $-3.6$, $-2.16$, $-1.296\,\%$.

\textbf{14.} $P = 200 - Y = 100 + z/2$: 110, 106, 103.6, 102.16, 101.296.

\textbf{15.} $10 + 6 + 3.6 + 2.16 + 1.296 = 23.056$, about 23\,\% of a year's \omterm{def:g12:aggregate-demand-and-supply:potential-output}{potential output}.

\textbf{16.} $z_2 = 20 \times 0.2 = 4$: gap $-2\,\%$.

\textbf{17.} Faster adjustment shrinks the loss: $10 + 2 + 0.4 + 0.08 + 0.016 \approx 12.5\,\%$ instead of 23\,\%.

\textbf{18.} Short run: $Y = 180 - Y$, $Y = 90$, $P = 90$. Long run: $Y = 100$, $P = 80$.

\textbf{19.} The demand shock lowers prices, first to 90 and then to 80; the supply shock raises them, to 110, and they fall back towards 100 as costs adjust.

\textbf{20.} The \omterm{def:g12:aggregate-demand-and-supply:potential-output}{output gap} is $-10$, $-6$, $-3.6$, $-2.16$ and $-1.3\,\%$ in years 1 to 5, and the cumulative output lost is about 23\,\% of a year's \omterm{def:g12:aggregate-demand-and-supply:potential-output}{potential output}.
\end{solution}
